% ext-lattice-kernel -- cited input, carried as an explicit hypothesis (state FROZEN). Not proved here.
% Source: local-times.tex:5-19 (flat 345-359), display eq:kernel-asymptotics.
\paragraph{Paper (verbatim).} ``Let $P$ be the simple random walk transition operator on $\Z^d$.
For $d=2$, let $b$ be the potential kernel normalized by $b(0)=0$. For $d\geq3$, let $b=-G$, where
$G(x)=\sum_{j\geq0}P^j(0,x)$ is the Green function. In both cases $(P-I)b=\ind_{\{0\}}$. The lattice
potential estimates of [Lawler--Limic, Theorem~4.3.1, Corollary~4.3.3, and Theorem~4.4.4] give, as
$|x|\to\infty$, $b(x)=\frac2\pi\log|x|+\kappa+O(|x|^{-2})$ for $d=2$, and
$b(x)=-\frac{2}{(d-2)\omega_d}|x|^{2-d}+O_d(|x|^{-d})$ for $d\geq3$, where $\kappa$ is the planar
potential-kernel constant.''
\paragraph{What is assumed.} For $d=2$: the partial sums $\sum_{j<M}[P^j(0,0)-P^j(0,x)]$ converge
for every $x$ to some $b(x)$, and $|b(x)-\frac2\pi\log|x|-\kappa|\le C|x|^{-2}$ for $|x|\ge R$. For
$d\ge3$: $|G(x)-\frac2{(d-2)\omega_d}|x|^{2-d}|\le C|x|^{-d}$ for $|x|\ge R$.
\paragraph{Truth check (vacuity).} Lawler--Limic Theorem 4.3.1 for simple random walk gives
$G(x)=\frac{d\Gamma(d/2)}{(d-2)\pi^{d/2}}|x|^{2-d}+O(|x|^{-d})$. Since
$\omega_d=\pi^{d/2}/\Gamma(d/2+1)$, this constant equals $2/((d-2)\omega_d)$; at $d=3$ both give
$3/(2\pi)$. Lawler--Limic (4.15) defines $a(x)=\lim_N[\sum_{n\le N}p_n(0)-\sum_{n\le N}p_n(x)]$. For
bipartite walks and odd $x$ the series converges only as a limit of partial sums, so the External
uses the limit of partial sums, not absolute summability. Theorem 4.4.4 gives, for simple random
walk, $a(x)=\frac2\pi\log|x|+\frac{2\gamma+\log8}\pi+O(|x|^{-2})$. Both clauses are therefore true, and
$\kappa$ exists. For $d\ge3$ summability of $P^j(0,x)$ is proved in the library
(\texttt{summable\_srwHeat}), so it is not assumed.
\paragraph{Transcription notes.} $b=-G$ turns the $d\ge3$ line into a statement about $G$; the
absolute value makes the sign immaterial. $1\le R$ encodes ``as $|x|\to\infty$''. $(P-I)b=\ind_{\{0\}}$
is derived, not assumed: for $d\ge3$ from \texttt{walkOp\_srwGreenInf}, and for $d=2$ by telescoping the
partial sums, which is node \texttt{s-kernel-props}. Corollary 4.3.3 (gradients) is not assumed:
eq:gradient is derived from the External, which is node \texttt{s-gradient}.
